\subsection{A validation protocol, and the objection it has to answer}
\label{sec:validation}
%==========================================================================

\emph{This is a protocol, not a completed study. Nothing in it has been run.}

\subsubsection{The obvious objection}
\label{sec:why_not_breath}

A reader who has followed Section~\ref{sec:spec} to its last row should ask the obvious
question: \textbf{if a breath test measures oxidative capacity directly, why identify
anything from records at all?}

The question is fair and it has three answers, in decreasing order of how much they settle.

\paragraph{Capacity is not margin.} The breath test measures how fast a patient oxidises a
labelled bolus. What the dosing decision turns on is whether that capacity exceeds what the
patient's \emph{own} load demands, which is $\mu = \theta - \thc$ and not $\theta$. The
threshold carries $\spk$ and the tolerable level, or in the equilibrium reduction the fold
of the patient's own clearance curve; either way it is a property of the patient's load and
diet, which the test does not see. A test result of ``forty per cent of normal oxidation''
does not say whether this patient, on this protein intake, has margin.

\paragraph{The test is not where the patients are.} The cohort a dosing decision has to be
made in is the cohort that exists: children on chronic management, in clinics that measure
acylcarnitines and organic acids and do not perform stable-isotope breath tests. A method
that requires an instrument nobody has is a method for a later cohort. The records exist
now, which is the whole reason this note is written on carnitine.

\paragraph{Scale transfer is not established.} The published test is in methylmalonic
acidemia~\cite{breath_test_mma}, % C48
where the block is one step further down the same pathway. Whether its reading maps onto
$\theta$ as this note defines it --- a maximal velocity at the plasma level, in a marker's
units --- is exactly the transfer question Remark~\ref{rem:no_km_orig} refuses to assume
elsewhere, and it would be inconsistent to assume it here.

\textbf{None of this makes the test dispensable.} It makes it the reference standard rather
than the replacement, which is what a validation protocol needs.

\subsubsection{The protocol}

\paragraph{Design.} Paired, within patient, single arm. Each participant supplies (i) a
marker-based estimate of $\theta$ from a record meeting Table~\ref{tab:spec}, and (ii) a
$1$-$^{13}$C-propionate breath test at the same visit. The comparison is of the two
estimates on a common scale, after the monotone rescaling of Section~\ref{sec:invariance},
so what is being validated is \emph{agreement up to a consistent distortion}, which is all
the decision requires.

\paragraph{Data per patient.} Exactly Table~\ref{tab:spec}: the marker on the specified
sampling schedule through at least one excursion, the load recorded to the stated precision
with enteral and parenteral sources separated, and one breath test at entry with a repeat at
any dose change.

\paragraph{Pre-specified agreement threshold.} The comparison is declared successful if the
limits of agreement on the log scale lie inside $\pm \ln 2$, that is if the two methods
agree within a factor of two on $\theta$. That is a weak threshold by the standards of assay
validation and a demanding one here, because Remark~\ref{rem:fibre} says a perfectly
resolved band still leaves $\theta$ open across a factor given by \eqref{eq:span}, which at
the illustrative values is $169$.

\paragraph{Power, against the reference standard's own precision.} The breath test's
repeatability is a coefficient of variation of $6.9\%$ in controls and $12.8\%$ in
patients~\cite{breath_test_mma}. % C48
A validation cannot resolve agreement more finely than the reference standard resolves
itself, so $12.8\%$ is a floor on the width of any limits of agreement, whatever the marker
side contributes. Two consequences follow and both should be stated before any data are
collected. First, a study powered to detect disagreement smaller than about a quarter
(roughly two reference standard deviations) is not powered against anything. Second, the
marker side's contribution is the quantity of Section~\ref{sec:collapse}, and at the
accessible designs it is an order of magnitude larger than the reference standard's, so the
study would be measuring the marker method's imprecision and not the agreement between two
comparable methods.

\paragraph{What would make the protocol worth running --- and it is no longer a
decompensation cohort.} An earlier version of this paragraph asked for a cohort observed
through decompensations with the load recorded throughout, on the grounds that only such a
cohort meets the excursion requirement of Table~\ref{tab:excursion}. Two findings retire
that: the excursion is not observed at the magnitude required, and the model loses its
stable equilibrium over most of the band in question (Section~\ref{sec:spec}). What
replaces it is Tier~1 of Table~\ref{tab:spec}: a same-day paired repeatability arm, which
is smaller than a natural history study by every measure that matters.

\paragraph{And the cohort arithmetic is what settles it.} Placing the fraction of patients
with $\mu<0$ within ten percentage points requires correcting for per-patient
misclassification, which inflates the variance by $(1-2p)^{-2}$. At the therapy-channel
design's simulated sign error near the boundary the requirement is about $2\,200$ patients,
against roughly $4\,900$ for the design it replaces. \textbf{Neither number is reachable and
the improvement does not change that.} The phase 1/2 trial enrolled
sixteen~\cite{trial_phase12}, % C16
and even in the highest-reported population this note cites, where newborn screening gives
an incidence near $1{:}14\,000$~\cite{nbs_saudi_national}, % C10
a cohort of two thousand is not a recruitment problem but an arithmetic one: it is a
substantial fraction of every diagnosed patient alive. In populations where the incidence is
lower --- which is most of them --- the arithmetic is worse, and this note does not quote a
figure for them because it has not verified one.

That is not an argument against the specification. It is an argument about what the
specification is \emph{for}. \textbf{A design that identifies $\theta$ per patient does not
need a cohort to decide that patient's dosing}; a cohort is needed only to estimate what
fraction of a population sits below the threshold, which is an epidemiological question and
not the clinical one this note opened with. The specification of Section~\ref{sec:spec} is
written for the first and should not be read as a proposal for the second.

\begin{remark}[A gap in the literature that the protocol depends on]
\label{rem:no_pa_breath}
\textbf{No published propionic acidemia breath-test cohort has been located.} The cohort
cited here is methylmalonic acidemia~\cite{breath_test_mma}, % C48
and the protocol above is written as though an equivalent series existed in PA. Confirming
whether one does, and if not what it would take to establish the test in this disease, is a
prerequisite rather than a detail. Until it is settled, the reference standard in the
protocol is an assumed one, and the $12.8\%$ quoted as its precision is transferred from an
adjacent disorder on the same pathway.
\end{remark}

\vspace{0.5em}
\hrule
\vspace{1em}

%==========================================================================
